We distinguish the state-ranking target from adjacent-centre change detection.
For each edge, a slot is positive if true support switches between consecutive
window centres; its score is the absolute change in emitted edge score.
Positive-change AP ranks strictly positive changes with tied levels kept
together and zero changes abstaining. Boundary-slot F1 uses a single
label-free threshold per fit, the $95$th percentile of raw nonself score
changes, shared by all methods. Stable-edge exceedances use that same
amplitude reference. These are stride-resolution diagnostics, not estimates
of sample-level boundary error; their exact coverage and conventions are in
Supplement Section~S17.2.

Figure~\ref{fig:localisation} shows that the metrics can disagree.
In Study A, forced-hard boundary-AP change is $-0.230$
$[-0.296,-0.162]$ at $D=300$ and $+0.184$ $[0.096,0.268]$
at $D=600$. The latter boundary-F1 change is only $+0.021$
$[-0.005,0.047]$ under the fixed amplitude reference.
In Study B at $D=600$, AP changes are $+0.068$, $+0.074$ and
$-0.022$ for fixed, random and irregular schedules; all three pointwise
intervals cross zero. Boundary F1 decreases by $0.077$, $0.068$ and
$0.087$, with all three intervals below zero. The positive state-AUROC
effects therefore do not establish improved boundary detection. Study C's
central $300/150$ grid likewise gives boundary-AP change $-0.055$
$[-0.132,0.026]$ and F1 change $-0.100$ $[-0.121,-0.080]$ at $D=600$,
despite its positive state-discrimination gain.

Hard emission reduces stable-edge threshold exceedances, but it also
suppresses genuine changes under this amplitude rule. Shrinkage preserves
change-ranking AP on these saved outputs while reducing change magnitudes;
it need not preserve threshold-based F1. The complete fixed-weight results,
stable-present/stable-absent denominators and all C grids are retained in
the numerical supplement. No absolute AP comparison is made across distinct
stride grids: effects compare each method with raw on the same slots.
